\documentclass{ltxdoc}

\usepackage{csquotes}
\usepackage{xparse}
\usepackage{xcolor}
\usepackage{tikz}
\usepackage{wrapfig}

\EnableCrossrefs

\CodelineIndex
\RecordChanges

\newcommand{\defaultValue}[1]{\textcolor{gray}{\textit{(#1)}}}
\DeclareDocumentCommand{\cmdArg}{s m}{\textcolor{gray}{\IfBooleanF{#1}{\{}\textit{#2}\IfBooleanF{#1}{\}}}}
\DeclareDocumentCommand{\confOption}{m v}{\noindent\textsf{\textbf{#1}}\IfValueT{#2}{\hfill\textit{\textcolor{gray}{default: #2}}}\\}

\setlength{\parskip}{1ex}
\setlength{\parindent}{0mm}


\changes{v1.0}{2017/08/28}{Initial version}



\title{The \textsf{RVTM} class}
\author{By Harvey Sheppard \\ \texttt{harveysheppard75 [at] yandex [dot] com}}


\begin{document}
	
	\maketitle
	
	\begin{abstract}
		The \textsf{RVTM} class is a class with some sort of modern style that I wrote for my brother.
	\end{abstract}
	
	\tableofcontents
	
	
	\section{How to Use the Configuration Options}
	The configuration options are created using the amazing \textsf{pgfkeys} system (part of \textsf{tikz}).
	
	To change the value of one configuration option, follow the following syntax:
	\begin{verbatim}
		\pgfkeys{/configuration/option = value, /configuration/option2 = value}
	\end{verbatim}
	
	
	
	\section{Global Ooptions}
	
	\confOption{/rvtm/language}|english|
	changes the language of the document.
	You can also change the language of the document using \verb|\setdefaultlanguage{<language>}|.
	
	\section{Layout Commands}
	
	\begin{macro}{\smallmarginlayout}
		Triggers the samll margin layout.
		Use \verb|\restoregeometry| to come back to the normal layout.
	\end{macro}
	
	
	
	\section{Font Commands}
	
	\begin{macro}{\geometricslim}
		Change the font to a geometrical slim font.
	\end{macro}
	
	\begin{macro}{\geometricbold}
		Change the font to a geometrical bold font.
	\end{macro}
	
	\begin{macro}{\highcontrastfont}
		Change the font to a font with big differences between the full and the loose parts of the letters.
	\end{macro}
	
	\begin{macro}{\rvtmtocchapterpagefont}
		This command defines the font of the numbers in the table of content.
	\end{macro}
	
	
	
	\section{Headers and Footers}
	
	\confOption{rvtm/footer/page number shape/line width}|1pt|
	The width of the line of the shape around the page number.
	
	
	
	\section{Chapters}
	
	\confOption{/rvtm/chapters/fontsize}|1.5cm|
	
	\confOption{/rvtm/chapters/line length}|\linewidth|
	The length of the horizontal lines above and below the chapter title.
	
	\confOption{/rvtm/chapters/line width}|2pt|
	The width of the horizontl lines above and below of the chapter text.
	
	\confOption{/rvtm/chapters/space between text and line}|2cm|
	
	
	
	\section{Epigraph}
	
	\confOption{/rvtm/epigraph/width}|.75\linewidth|
	
	\begin{macro}{\epigraphrule}
		Change this length to change the length of the epigraph rule's width. Default is 0pt.
	\end{macro}
	
	\begin{macro}{\rvtmepigraph}
		This command is used to typeset an epigraph after the chapters.
		\begin{description}
			\item[First Mandatory Argument] The quotation.
			\item[Second Optional Argument] The author of the quotation.
		\end{description}
	\end{macro}
	
	
	
	\section{Table of Content}
	
	
	\subsection{Left Column Configuration}
	
	\confOption{/rvtm/toc/left column width}|.5\textwidth|
	The width of the column containing the chapter's text.
	
	\confOption{/rvtm/toc/right column widh}|.45\paperwidth|
	The width of the column containing the chapter's page number.
	
	\confOption{/rvtm/toc/title/fontsize}|1.5cm|
	The font size of the title of the table of content (left column).
	
	\confOption{/rvtm/toc/title/font}|\geometricslim\addfontfeature{LetterSpace=-10.0}|
	The code used to configure the font of the table of content's title (left column).
	
	\confOption{/rvtm/toc/title/translation string}|\GetTranslation{Overview}|
	The string to use for the table of content's title.
	The default value makes use of the default strings defined by the \textsf{translations} package.
	
	\confOption{/rvtm/toc/title/bar thickness}|3mm|
	The thickness of the rules at the top and bottom of the table of content.
	
	\confOption{/rvtm/toc/title/bar length}|.8\linewidth|
	The length of the lines at the top and bottom of the table of content.
	A relative length can be defined because the table of content is put inside a minipage whose width is defined by \textsf{/rvtm/toc/left column width}.
	
	\confOption{/rvtm/toc/chapter label/fontsize}|5mm|
	
	\confOption{/rvtm/toc/chapter label/color}|darkGrey|
	See the \textsf{yMainColor} package for the definition of this color.
	
	\confOption{/rvtm/toc/chapter label/font}|\itshape|
	
	\confOption{/rvtm/toc/chapter text/fontsize}|5mm|
	
	\confOption{/rvtm/toc/chapter text/color}|middleGrey|
	See the \textsf{yMainColor} package for the definition of this color.
	
	\confOption{/rvtm/toc/chapter text/font}|\geometricbold\bfseries|
	
	
	\subsection{General Commands}
	
	\begin{macro}{\rvtmtocbar}
		Prints a bar as the ones found at the top and bottom of the table of content.
	\end{macro}
	
	\begin{macro}{\rvtmtoc}
		You must use this command to print the table of content and not the normal \LaTeX\ command.
	\end{macro}
	
	
	\subsection{Right Column Configuration}
	
	\confOption{/rvtm/toc/chapter page/color}|rvtmtoccolor|
	The color used for the page number.
	
	\confOption{/rvtm/toc/chapter page/yshift}|10mm|
	The maximum vertical overlap between two consecutive page number.
	This value is then randomized for each chapter page number.
	
	\confOption{/rvtm/toc/chapter page/xshift}|5mm|
	The maximum value of the horizontal shift a page number might receive.
	This value is then randomized for each chapter page number.
	
	\confOption{/rvtm/toc/chapter page/initial yshift}|7mm|
	How much above the page the toc should begin.
	
	\confOption{/rvtm/toc/chapter page/line width}|2pt|
	
	\confOption{/rvtm/toc/chapter page/distance from chapter text to line}|5mm|
	
	\confOption{/rvtm/toc/chapter page/font}|\rvtmtocchapterpagefont\addfontfeatures{LetterSpace=-15.0}|
	The code to customize the font of the chapter page number.
	
	It is possible to define \textit{by hand} all the lengths required to typeset the page numbers in the right column.
	All the length described after receive default values just after the \verb|\begin{document}|.
	Most of those values are randomized length of a basis length.
	Also note that the numbers are alternatively aligned to the left and to the right beginning on the left.
	If you want to redefine those length, please do it after \verb|\begin{document}|.
	It is possible to redefine only a few of the length described after.
	The complete list of required length is the following:
	\begin{description}
		\item[\verb|\@rvtm@toc@chapter@page@<x>@fontsize|] The font size of the page number of the chapter \verb|<x>|.
		\item[\verb|\@rvtm@toc@chapter@page@<x>@xshift|] The horizontal shift of the page number of the chapter \verb|<x>|.
		\item[\verb|\@rvtm@toc@chapter@page@<x>@yshift|] The vertical shift of the page number of the chapter \verb|<x>|.
		\item[\verb|@rvtm@toc@chapter@page@<x>@line@end@xshift|] (! This is not a length but a command !) The percentage of the current chapter page number font size that the end of the line should overlap over the chapter page number.
		\item[\verb|@rvtm@toc@chapter@page@<x>@line@end@angle|] (! This is not a length but a command !) the angle in degrees, centered at the chapter page number center, that shouold be used to indicate where the line will arrive in the chapter page number.
		
		\begin{tikzpicture}
			\node [draw, dashed, font=\fontsize{1cm}{1.2cm}\selectfont] (text) {\textcolor{gray}{X}};
			\draw (text.center) -- (text.east);
			\draw (text.center) -- (text.150);
			\draw ([xshift=-2mm, yshift=1.1mm]text.center) arc (150:0:2mm);
			\fill [red] (text.150) circle (.5mm);
		\end{tikzpicture}
		
		In the drawing above, the grey \enquote{X} represents the chapter page number and the red dot indicates where the line joining the chapter text and the chapter page number will end.
		
		Please note that if you want to define custom entry angle, you have to add \verb|\toggletrue{rvtmtoccustomchapterpagelineentryangle}| after the \verb|\begin{document}|.
	\end{description}
	
	
	
	\section{Special Pages}
	
	
	\subsection{Framed Page}
	
	\confOption{/rvtm/framedpage/line width}|2mm|
	The width of the framing line.
	
	\confOption{/rvtm/framedpage/line color}|black|
	
	\confOption{/rvtm/framedpage/text width}|.5\textwidth|
	
	\confOption{/rvtm/framedpage/vertical offset}|2cm|
	The vertical distance between the page border and the framing line.
	
	\confOption{/rvtm/framedpage/horizontal offset}|2cm|
	The horizontal distance between the page border and the framing line.
	
	\confOption{/rvtm/framedpage/offset}|2cm|
	A shortcut to define in one call both the vertical and horizontal offset.
	
	\confOption{/rvtm/framedpage/above space}|\vfill|
	The vertical space added before the actual content of the framed page.
	
	\confOption{/rvtm/framedpage/left space}|\hfill|
	The horizontal space added before the actual content of the framed page.
	
	\confOption{/rvtm/framedpage/right space}|\hfill|
	The horizontal space added after the actual content of the framed page.
	
	\confOption{/rvtm/framedpage/below space}|\vfill|
	The vertical space added after the actual content of the framed page.
	
	\begin{environment}{framedpage}
		This environment creates a page with a frame.
		The text inside the environment will be typeset in a minipage.
	\end{environment}
	
	\confOption{/rvtm/framedpage/title/font}|\geometricbold|
	The code applied to the framed page's title font.
	
	\confOption{/rvtm/framedpage/title/size}|2em|
	
	\confOption{/rvtm/framedpage/title/below skip}|2mm|
	
	\begin{macro}{\framedpagetitle}
		This command typeset a nice title for the framed pages.
		Feel free to create your own titles styles.
		This is only a proposal.
		\begin{description}
			\item[First Mandatory Argument] The title's text.
		\end{description}
	\end{macro}
	
	
	\subsection{Black Pages}
	
	\confOption{/rvtm/blackpage/background}|black|
	The background color.
	
	\confOption{/rvtm/blackpage/text color}|white|
	
	\confOption{/rvtm/blackpage/text width}|.5\textwidth|
	
	\confOption{/rvtm/blackpage/above space}|\vfill\vfill|
	The space above the text.
	
	\confOption{/rvtm/blackpage/below space}|\vfill|
	The space below the text.
	
	\confOption{/rvtm/blackpage/left space}|\strictpagechecktrue\checkoddpage\ifoddpage\relax\else\hfill\fi|
	Horizontal space to add before the text.
	The default value aligns the text to the inner of the page.
	
	\confOption{/rvtm/blackpage/right space}||
	
	\confOption{/rvtm/blackpage/vertical offset}|0mm|
	The vertical offset between the background color and the border of the page.
	
	\confOption{/rvtm/blackpage/horizontal offset}|0mm|
	The horizontal offset between the background color and the border of the page.
	
	\confOption{/rvtm/blackpage/offset}|0mm|
	A shortcut to define both the vertical and horizontal offset.
	
	\begin{environment}{blackpage}
		This environment add a colored background to a new page and print the text inside the environment inside a minipage.
	\end{environment}
	
	
	
	\section{Backgroud Photo for Double Page}
	
	\confOption{/rvtm/bpdp/photo width}|2\paperwidth|
	
	\confOption{/rvtm/bpdp/photo height}|\paperheight|
	
	\confOption{/rvtm/bpdp/page style}|empty|
	Used to change the header and footer type.
	If you want to use custom ones, please refer to the \textsf{fancyhdr} package.
	
	\begin{environment}{bpdp}
		Prints a photo on the next double page.
		Half of the photo will be on the left page, the other half will be on the right page.
		
		Please note that you must use \verb|\bpdpnewpage| exactly once inside the \textsf{bpdp} environment.
		This command begins the second backgrounded page.
		
		\begin{description}
			\item[First Mandatory Argument] The path to the photo to display.
		\end{description}
	\end{environment}
	
	\begin{macro}{\bpdpnewpage}
		Starts the second backgrounded page.
		Must appear exactly once inside the \textsf{bpdp} environment where you would normally put a \verb|\newpage|.
	\end{macro}
	
	
	
	\section{Horizontal Dotted Lines}
	
	\confOption{/rvtm/dotted line/fill}|black|
	The color of the dots.
	
	\confOption{/rvtm/dotted line/radius}|1mm|
	The radius of the dots.
	
	\confOption{/rvtm/dotted line/space between dots}|2mm|
	
	\confOption{/rvtm/dotted line/above skip}|\baselineskip|
	The skip above the dotted line.
	
	\confOption{/rvtm/dotted line/below skip}|\baselineskip|
	The sapce below the dotted line.
	
	\begin{macro}{\hdottedline}
		Fill the rest of the line with dots.
		This command is to be used like \verb|\hrule|.
		
		\begin{description}
			\item[First Optional Argument] A style for an horizontal dotted line.
			\item[Second Optional Argument] Some configuration options (as described just above) for the horizontal dotted line, where the \verb|/rvtm/dotted line/| prefix was removed.
		\end{description}
		
		Note: you can define custom dotted line style using \begin{verbatim}
		\rvtmnewstyle{<style name>}{
			dotted line/.cd,
			<configuration option for dotted line without prefix> = value,
		}
		\end{verbatim}
		
		The \textsf{white dotted line} style defines a line with white dots.
	\end{macro}
	
	
	
	\PrintChanges
	\PrintIndex
\end{document}